There are different properties that we should consider when we are performing a SAT encoding. We have already claimed some of them, as the number of 
variables and the total cardinality of clauses. However, there is something more that we have to keep in mind while we are doing so, that is 
\textbf{constraint propagation}. \vspace{7pt}

Let's consider an encoding $E$ of a constraint $C$, such that there is a correspondence between the assignments of the variables in $C$ with boolean
assignments of the variables in $E$. $E$ is arc-consistent if:
\begin{itemize}
    \renewcommand{\labelitemi}{-}
    \item Whenever a partial assignment is inconsistent wrt $C$ (cannot be extended to a solution of $C$), unit propagation in $E$ causes conflict.
    \item Otherwise, unit propagation in $E$ discards arc-inconsistent values.
\end{itemize}
\begin{example}
    e.g. Arc-Consistency with $\text{AtMostOne}([x_1, x_2, \dots, x_n])$ constraint. \vspace{7pt}

    $\text{AtMostOne}([x_1, x_2, \dots, x_n])$ is arc-consistent when:
    \begin{itemize}
        \renewcommand{\labelitemi}{-}
        \item If there are two variables $x_i$ and $x_j$ assigned to $T$ then unit propagation should give a conflict.
        \item If there is one $x_i$ assigned to $T$ then all the other $x_j$ should be assigned to $F$ by unit propagation.
    \end{itemize} \vspace{7pt}

    Note: everything must be done by unit propagation, without any additional operation from the SAT solvers.
\end{example} \vspace{3.5pt}

\textbf{Are the AtMostOne(\dots) encodings we have seen so far arc-consistent encodings?} \\ 
All of them are arc-consistent except one: \textbf{bitwise encoding}. Now we are gonna to see an example to state the inconsistency of this encoding, considering
always unit propagation as the main operation.